\section{Contenido de la memoria}
\label{Contenido de la memoria}

    La presente memoria se encuentra estructurada en diversos capítulos que recogen, de manera sistemática y progresiva, las fases de análisis, diseño, desarrollo y validación de Colivi, una plataforma web y móvil asistida por inteligencia artificial orientada a la gestión integral del alquiler en régimen de coliving y a la convivencia entre residentes. A continuación, se resume de forma sucinta el propósito y contenido de cada una de las partes que integran este trabajo:

- Introducción: se contextualiza la creciente demanda del modelo residencial de coliving entre estudiantes y jóvenes profesionales, exponiendo las problemáticas actuales asociadas a la dispersión de herramientas digitales durante la búsqueda de piso y la convivencia diaria.

- Justificación: se detallan las razones técnicas y sociales que motivan la elección de esta propuesta, destacando la oportunidad de resolver la fragmentación operativa del alquiler compartido unificando el ciclo completo en una sola solución con soporte inteligente.

- Objetivos: se formula el objetivo general del proyecto y se desglosan los objetivos específicos, delimitando el alcance funcional de la plataforma tanto en la vertiente web como en su empaquetado móvil.

- Planificación: se describe la organización temporal del desarrollo a partir de un ciclo de vida clásico de ingeniería del software, exponiendo las fases del proyecto, el cronograma de trabajo planificado en 300 horas lectivas y la trazabilidad de las dependencias lógicas.

- Aplicaciones similares: se ofrece un análisis comparativo frente a soluciones existentes en el mercado (portales de búsqueda inmobiliaria, plataformas de alojamiento temporal y utilidades de reparto de gastos), delimitando los criterios funcionales que diferencian y fundamentan la propuesta de Colivi.

- Herramientas y tecnologías: se justifican las decisiones metodológicas y la selección del stack tecnológico empleado, profundizando en el backend con Java 21 y Spring Boot 3, el frontend reactivo con React, TypeScript y Tailwind CSS, la persistencia relacional con PostgreSQL, la integración del protocolo Model Context Protocol (MCP) con Spring AI y Groq, y la adaptación móvil mediante Capacitor.
 
- Análisis y diseño del sistema: se documentan los diferentes roles de usuario, el catálogo riguroso de requisitos funcionales y no funcionales, los diagramas de casos de uso, el modelo entidad-relación de la base de datos y la arquitectura distribuida del sistema desacoplada bajo principios SOLID.

- Implementación del sistema: se expone detalladamente el desarrollo técnico modular de la plataforma, abordando el control de acceso sin estado con JWT, el catálogo con mapas interactivos, el flujo de reservas con fianza integrada y control de concurrencia, los módulos de convivencia del hogar (gastos compartidos, tareas domésticas con asignación rotativa y mensajería interna) y el microservicio MCP que da soporte al asistente conversacional.

- Adaptabilidad móvil y despliegue: se describe la optimización ergonómica de la interfaz de usuario para dispositivos móviles, el empaquetado del aplicativo para Android y la configuración de la infraestructura mediante contenedores Docker y servicios cloud en Microsoft Azure con automatización CI/CD.

- Plan de pruebas y aseguramiento de la calidad: se detalla la estrategia de verificación adoptada, cubriendo las pruebas unitarias y de integración de la lógica de negocio, las comprobaciones de los endpoints seguros y el análisis de incidencias técnicas resueltas durante el desarrollo.

- Estudio de costes: se desglosa el presupuesto global estimado para la ejecución del proyecto, cuantificando los costes derivados del esfuerzo humano (horas de ingeniería), la infraestructura cloud, las herramientas de desarrollo y los recursos materiales requeridos.

- Conclusiones y líneas futuras: se sintetizan las metas alcanzadas y el grado de cumplimiento de los objetivos planteados, proponiendo posibles vías de evolución técnica para extender las capacidades operativas y de inteligencia del sistema.

- Bibliografía: se recopilan formalmente todas las referencias bibliográficas, documentación técnica oficial, especificaciones de protocolos, estándares y fuentes académicas consultadas a lo largo de la investigación y construcción de la memoria.